% ============================================================
% M3 S2 导学件·波1 臂6 片件 —— 课时10 曲线与方程（2.4）＝批C 四片之一（母版照抄制）
% 生成：M3 成卷轮 S2 波1 臂6｜2026-09-14｜编译 cwd＝本目录（中文路径禁 \input 绝对路径）
%   xelatex main-true.tex ×2 → 含详解印本；xelatex main-false.tex ×2 → 纯题版（单源双本）
% 输入链（全只读）：题面库 课时10-2.4曲线与方程.md（题面逐字装配）＋ -答案侧.md（值逐字回嵌）
%   ＋ 定稿/批C-课时10-2.4曲线与方程.md（详解回嵌源，§9 全文区）；toolchain＝qp-m3.sty（钉后版
%   md5 7c3930362be8a0a2bdf21bbf8ac16573，本地挂载副本，破损即停）。
% 母版体例：详见 成卷/导学件/课时01-坐标法/母版体例说明.md（结构骨架/宏用法/门谱跑法/常见坑）。
% 键账：件内 ansblock 键集 ≡ 题面库片键集 ≡ manifest 键序（21 键，零缺漏零多余）；
%   预习位零键零答案（口令④裁定前口径）；印面号 1—21＝探究 1—16＋课堂评价 17—21，
%   键↔号映射见 值台账-课时10.json。
% 括线判据（承重墙禁放宽）：估高＞8 行（chars/23 栏宽口径）→ 括线模，逐键读数入值台账；
%   其余灰底。本片括线键读数见值台账（门谱实测）。
% 观察注落实（军师令·补产批CA 核账 §C.4 随片下发）：G2＝装配层加「定义提示卡」（\\zhuzhu，
%   两档等形不泄答）在印；拓2/拓3/拓8/拓10/中7＝非本件键域（拓展册席/置换移出席），
%   观察注随档移交拓展册工位，本件落实数 0；逐条读数见 值台账-课时10.json「观察注落实」。
% 门谱读数：见过程件 工作区/_tmpM3S2波1臂6_0914/（编译 log、门谱输出、键表）。
% ============================================================
\documentclass[fontset=none]{ctexart}
\usepackage{qp-m3}
% —— 印面逐字符号路由补钉（探针实证 0914：书宋缺 U+2212，▱ 诸字库皆缺→NSC 子块；
%     禁改 sty 本体保 md5；无 ▱/− 用件的片可整块照抄，零副作用） ——
\xeCJKDeclareCharClass{Default}{"2212}%
\setCJKfamilyfont{nscblk}[Path=C:/提示词/工作区/字替对照-0909/variantF/fonts/,BoldFont=NSC-w600.ttf]{NSC-w500.ttf}%
\xeCJKDeclareSubCJKBlock{nscblk}{"25B1}%
\newcommand{\pxparallelogram}{{\CJKfamily{nscblk}▱}}%
% —— 断行微调（纯题档/详解档三零门：CJK 胶收缩 0.01→0.05em；应急伸展曾试 12em/2em，
%     2026-09-14 实测撤行回落 sty 默认 1em 三零仍守，故不设该行——补丁最少化；
%     只动伸缩分量不动自然宽度，版面纹理零漂；Underfull 治理读数见过程件） ——
\xeCJKsetup{CJKglue={\hskip 0.10em plus 0.08em minus 0.05em}}%
% —— 页脚件名（qp-headfoot 复用入口） ——
\renewcommand{\qpzhangming}{第二章\quad 平面解析几何}%
\renewcommand{\qpjianming}{导学件}%

\begin{document}

\zhangtitle{第二章\quad 平面解析几何}
\jietitle{2.4 曲线与方程}
\mubiaoline
\mubiaomu{1}{理解曲线的方程、方程的曲线的概念，能判断点与方程（曲线）的对应关系，会说清纯粹性与完备性；}
\mubiaomu{2}{掌握求动点轨迹方程的一般步骤（建系、设点、列等量关系、化简、检验），会用直接法、相关点法求轨迹方程并剔除不合条件的点；}
\mubiaomu{3}{会由方程研究曲线的对称性、范围与坐标轴交点，能对含参方程定形分类讨论，体会“几何问题代数化、代数结果几何化”的解析几何思想.}

\vspace{1.7mm}
\begin{multicols}{2}
\emergencystretch=1em

% ============================================================
% 课前预习（知识点导学填空；零键零答案——预习键位按检查口令④裁定前口径，
% \kongbai 空线＋\zhentib 判断题面形，两档等形不泄答）
% ============================================================
\huaxing{课}{前}{预}{习}{知识导学\quad 素养初识}

\zsd{一}{曲线的方程与方程的曲线}

\tiaomu{1}{一般地，在直角坐标系中，若曲线\(C\)上的点与二元方程\(f(x,y)=0\)的实数解满足：（1）曲线上的点的坐标都是这个方程的\kongbai{}；（2）以这个方程的\kongbai{}为坐标的点都在曲线上，那么这条曲线叫作方程的曲线，这个方程叫作曲线的方程．\zhuzhu{（1）说纯粹性（不杂），（2）说完备性（不漏）；两者缺一不可，判断对应关系时双向都查.}}

\zsd{二}{求动点轨迹方程的一般步骤}

\tiaomu{1}{建系（选基本量明显的坐标系）→设点（设动点坐标\((x,y)\)）→列式（把几何条件写成\kongbai{}）→化简（同解变形）→\kongbai{}（查纯粹性与完备性，除去不合条件的点）．\zhuzhu{斜率、比值含分母的，先写存在条件（如分母不为零、不共线）；与定点重合、三点共线常是剔点来源.}}

\tiaomu{2}{相关点法：所求点\(P\)随已知曲线\(F(x,y)=0\)上的动点\(Q\)而动时，先用\(P\)表示\(Q\)，代入\(Q\)满足的方程即得\(P\)的\kongbai{}方程．\zhuzhu{中点、对称、重心等问题先用中点公式（对称式）反解旧点坐标，再整体代入.}}

\zsd{三}{由方程研究曲线}

\tiaomu{1}{通过方程研究曲线的\kongbai{}（换\(x\to-x\)、\(y\to-y\)检验）、范围（由方程解出\(x\)、\(y\)的取值）与曲线和坐标轴（或直线）的\kongbai{}，进而认识曲线的形状．\zhuzhu{含绝对值、含参方程常需分段或配方，配方定形是本章后几节课的常用入口.}}

\zhenhead{判断正误(正确的打\gou{},错误的打\cha{})}

\zhentib{(1)}{以方程\(f(x,y)=0\)的解为坐标的点都在曲线\(C\)上，且曲线\(C\)上点的坐标都是\(f(x,y)=0\)的解，则\(C\)叫作方程\(f(x,y)=0\)的曲线．}

\zhentib{(2)}{方程\(y=\sqrt{x^{2}}\)与\(y=x\)表示同一条曲线．}

\zhentib{(3)}{到两坐标轴距离相等的点的轨迹方程是\(y=x\)．}

\liubai[9mm]{此处书写}

% ============================================================
% 课中探究（考点探究〔例+变式〕＝练/拓槽 16 键；题后紧跟答案＝ansblock 内嵌，
% 值照答案侧逐字，详解承定稿 §9 回嵌）
% ============================================================
\huaxing{课}{中}{探}{究}{考点探究\quad 素养小结}

% ---------- 探究点一 曲线的方程与方程的曲线 ----------
\tjdnr{一}{曲线的方程与方程的曲线}{例\textbf{1}}{简单(知识点一)}{判断\(P(\sqrt{3},0)\)，\(Q(-2,3)\)是否在方程\(4x^{2}+3y^{2}=12\)的曲线上．}
\xiexwei{7mm}

\begin{ansblock}[2章-练-课时10-简1]
% ans:2章-练-课时10-简1
\ansitem{1}{P在曲线上；Q不在曲线上}
\ansnote{详解}{曲线上的点的坐标必为方程的解，反之方程的解对应曲线上的点，故代入即判．\(P(\sqrt{3},0)\)：\(4\times(\sqrt{3})^{2}+3\times 0^{2}=12+0=12\)，是方程的解，\(P\)在曲线上．\(Q(-2,3)\)：\(4\times(-2)^{2}+3\times 3^{2}=16+27=43\neq 12\)，不是方程的解，\(Q\)不在曲线上．}
\end{ansblock}

\liB{变式\textbf{2}}{简单(知识点一)}{}{已知方程\((x+2)^{2}+(y-3)^{2}=r^{2}\)的曲线通过点\((-1,5)\)，求\(r^{2}\)的值．}
\xiexwei{7mm}

\begin{ansblock}[2章-练-课时10-简2]
% ans:2章-练-课时10-简2
\ansitem{2}{r²=5}
\ansnote{详解}{曲线通过点\((-1,5)\)，即\((-1,5)\)的坐标是方程的解：\((-1+2)^{2}+(5-3)^{2}=r^{2}\)，得\(r^{2}=1+4=5\)．}
\end{ansblock}

\liB{变式\textbf{3}}{简单(知识点一)}{}{到两坐标轴距离相等的点组成的轨迹的方程是\(x-y=0\)吗？为什么？}
\xiexwei{9mm}

{\ansblockgrayfalse
\begin{ansblock}[2章-练-课时10-简3]
% ans:2章-练-课时10-简3
\ansitem{3}{不是。x−y=0 只表示其中一条直线，轨迹应为 x−y=0 或 x+y=0（即|x|=|y|）；以 x−y=0 为轨迹方程不完备（漏了直线 x+y=0 上的点）}
\ansnote{详解}{设动点\(M(x,y)\)．到\(x\)轴距离为\(|y|\)、到\(y\)轴距离为\(|x|\)，条件\(|x|=|y|\iff x-y=0\)或\(x+y=0\)，轨迹是两条互相垂直的直线．方程\(x-y=0\)的解对应的点都在轨迹上（纯粹性成立），但轨迹上的点（如\((1,-1)\)）不都是该方程的解（完备性不成立）——它只描述第一、三象限角平分线，漏了第二、四象限角平分线．故“轨迹的方程是\(x-y=0\)”不对．}
\end{ansblock}}

\liB{变式\textbf{4}}{简单(知识点一)}{}{下列各组方程中，表示同一条曲线的是\nobreak(\kongwei)}
\bindp A．\(y=|x|\) 与 \(y=\sqrt{x^{2}}\)

\bindp B．\(y=x\) 与 \(y=\sqrt{x^{2}}\)

\bindp C．\(|y|=x\) 与 \(y=\sqrt{x^{2}}\)

\bindp D．\(x^{2}+y^{2}=0\) 与 \(y=0\)

\begin{ansblock}[2章-练-课时10-简7]
% ans:2章-练-课时10-简7
\ansitem{4}{A}
\ansnote{详解}{A：\(\sqrt{x^{2}}=|x|\)，两方程解集（图象）完全相同，表同一条曲线；B：\(y=x\)为直线，\(y=\sqrt{x^{2}}=|x|\)为折线（\(x<0\)时\(y=-x\)），异；C：\(|y|=x\iff x\geqslant 0\)且\(y=\pm x\)，为两条射线，而\(y=|x|\)还含\(x<0\)部分，异；D：\(x^{2}+y^{2}=0\iff\)点\((0,0)\)，\(y=0\)为整条\(x\)轴，异．故选：A．}
\end{ansblock}

\liB{变式\textbf{5}}{简单(知识点一)}{}{与方程 \(x^{2}+y^{2}=4\) 表示同一条曲线的是\nobreak(\kongwei)}
\bindp A．\(y=\sqrt{4-x^{2}}\)

\bindp B．\((x-2)^{2}+y^{2}=4\)

\bindp C．\(x^{2}+y^{2}-4=0\)

\bindp D．\(|x|+|y|=2\)

\begin{ansblock}[2章-练-课时10-简8]
% ans:2章-练-课时10-简8
\ansitem{5}{C}
\ansnote{详解}{C 与原方程仅差移项，同解变形，解集恒等，表同一条曲线；A：\(\sqrt{4-x^{2}}\) 要求 \(y\geqslant 0\)，仅为上半圆，缺 \(y<0\) 部分不完备，排除；B：圆心 \((2,0)\)、半径 2 的圆，与圆心在原点的圆非同一条曲线，排除；D：\(|x|+|y|=2\) 是以 \((2,0)\)、\((0,2)\)、\((-2,0)\)、\((0,-2)\) 为顶点的正方形边界，点 \((1,1)\) 在其上却不在圆上，排除．故选：C．}
\end{ansblock}

\liB{变式\textbf{6}}{简单(知识点一)}{}{判断下列命题真假并说明理由：}
\bindp (1) 方程 \(x^{2}+y^{2}=0\) 表示一条直线；

\bindp (2) \(|x|=|y|\) 的轨迹方程是 \(y^{2}=x^{2}\)；

\bindp (3) 若点 \(P(x_0,y_0)\) 的坐标是曲线 \(C\) 的方程 \(f(x,y)=0\) 的一组解，则 \(P\) 在 \(C\) 上．
\xiexwei{7mm}

\begin{ansblock}[2章-练-课时10-简10]
% ans:2章-练-课时10-简10
\ansitem{6}{(1) 假　(2) 真　(3) 真}
\ansnote{详解}{(1) \(x^{2}+y^{2}=0\) 当且仅当 \(x=y=0\)，表示原点一个点，非直线；(2) \(|x|=|y|\) 当且仅当 \(x^{2}=y^{2}\)（两侧均非负，平方同解），即 \(y^{2}=x^{2}\)，命题真；(3) 由曲线方程定义的完备性（以方程的解为坐标的点都在曲线上），命题真．}
\end{ansblock}

\xiaojie{①识别：判断点与方程（曲线）的对应、辨析“同一条曲线”、判断轨迹方程的纯粹性完备性，判归本题型.}
\bindp {\kaishu ②操作：代入即判点的进出；两方程比解集（比取值范围、比图象）而不是比外形；说理须双向——纯粹性（解入曲线）与完备性（点入解）逐一交代.}
\bindp {\kaishu ③收束：绝对值方程平方变形须两侧非负才同解；轨迹方程漏解（如漏另一条角平分线）是本类题最常见错误.}

% ---------- 探究点二 求轨迹方程·直接法 ----------
\tjdnr{二}{求轨迹方程·直接法}{例\textbf{7}}{简单(知识点二)}{设\(A(-1,-1)\)，\(B(3,7)\)，求线段\(AB\)的垂直平分线的方程．}
\xiexwei{9mm}

{\ansblockgrayfalse
\begin{ansblock}[2章-练-课时10-简4]
% ans:2章-练-课时10-简4
\ansitem{7}{x+2y−7=0}
\ansnote{详解}{设\(P(x,y)\)为线段\(AB\)垂直平分线上任意一点，则\(|PA|=|PB|\)：\((x+1)^{2}+(y+1)^{2}=(x-3)^{2}+(y-7)^{2}\)，展开：\(x^{2}+2x+1+y^{2}+2y+1=x^{2}-6x+9+y^{2}-14y+49\)，化简得 \(8x+16y-56=0\)，即 \(x+2y-7=0\)．反向检验：满足 \(x+2y-7=0\) 的点均有\(|PA|=|PB|\)，在\(AB\)的垂直平分线上（两性完备）．验算：\(AB\)中点\((1,3)\)：\(1+6-7=0\)成立；\(k_{AB}=(7+1)/(3+1)=2\)，所求直线斜率\(-1/2\)，\(2\times(-1/2)=-1\)成立．}
\end{ansblock}}

\liB{变式\textbf{8}}{简单(知识点二)}{}{求到两定点\(A(-1,2)\)，\(B(3,2)\)的距离之比为\(\sqrt{2}\)的点的轨迹方程．}
\xiexwei{12mm}

{\ansblockgrayfalse
\begin{ansblock}[2章-练-课时10-简5]
% ans:2章-练-课时10-简5
\ansitem{8}{(x−7)²+(y−2)²=32}
\ansnote{详解}{设\(M(x,y)\)，\(|MA|/|MB|=\sqrt{2}\)，即\(|MA|^{2}=2|MB|^{2}\)：\((x+1)^{2}+(y-2)^{2}=2[(x-3)^{2}+(y-2)^{2}]\)．展开：\(x^{2}+2x+1+(y-2)^{2}=2x^{2}-12x+18+2(y-2)^{2}\)，移项整理：\(x^{2}-14x+(y-2)^{2}+17=0\)，配方得 \((x-7)^{2}+(y-2)^{2}=32\)．轨迹为以\((7,2)\)为圆心、\(4\sqrt{2}\)为半径的圆（阿波罗尼斯圆雏形，此名可不向学生出现）．验算：圆与直线\(AB\)（\(y=2\)）交于 \(M_1(7-4\sqrt{2},2)\)、\(M_2(7+4\sqrt{2},2)\)．对\(M_1\)：\(|M_1A|=|7-4\sqrt{2}+1|=8-4\sqrt{2}=4(2-\sqrt{2})\)，\(|M_1B|=|4\sqrt{2}-4|=4(\sqrt{2}-1)\)（取绝对值），比值\(=(2-\sqrt{2})/(\sqrt{2}-1)=\sqrt{2}(\sqrt{2}-1)/(\sqrt{2}-1)=\sqrt{2}\)成立；对\(M_2\)同法比值亦\(\sqrt{2}\)成立．}
\end{ansblock}}

\liB{变式\textbf{9}}{简单(知识点二)}{}{已知\(\triangle ABC\)底边两端点\(B(0,6)\)、\(C(0,-6)\)，若这个三角形另外两边所在直线的斜率之积为\(-4/9\)，求点\(A\)的轨迹方程．}
\xiexwei{9mm}

\begin{ansblock}[2章-练-课时10-简6]
% ans:2章-练-课时10-简6
\ansitem{9}{x²/81+y²/36=1（x≠0）}
\ansnote{详解}{设\(A(x,y)\)，\(x\neq 0\)（\(x=0\) 时 \(A\) 在 \(y\) 轴上与 \(B\)、\(C\) 共线，不能成三角形且斜率不存在）．\(k_{AB}\cdot k_{AC}=[(y-6)/x]\cdot[(y+6)/x]=(y^{2}-36)/x^{2}=-4/9\)，去分母：\(9(y^{2}-36)=-4x^{2}\)，整理得 \(4x^{2}+9y^{2}=324\)，即 \(x^{2}/81+y^{2}/36=1\)（\(x\neq 0\)）．}
\end{ansblock}

\liB{变式\textbf{10}}{简单(知识点二)}{}{已知动点 \(M\) 到 \(x\) 轴的距离与到定点 \(A(0,4)\) 的距离相等，求 \(M\) 的轨迹方程。}
\xiexwei{9mm}

\begin{ansblock}[2章-练-课时10-简9]
% ans:2章-练-课时10-简9
\ansitem{10}{x²=8(y−2)}
\ansnote{详解}{设 \(M(x,y)\)．由题 \(|y|=\sqrt{x^{2}+(y-4)^{2}}\)．两边平方（两侧非负，等价）：\(y^{2}=x^{2}+y^{2}-8y+16\)，即 \(x^{2}=8(y-2)\)．检验：方程上任一点均满足定义（平方等价可逆），无杂点；\(A(0,4)\) 代入 \(16\neq 0\) 不在轨迹上，无需剔除．故轨迹方程为 \(x^{2}=8(y-2)\)．}
\end{ansblock}

\liB{变式\textbf{11}}{中档(知识点二)}{}{已知\(A(-2,0)\)、\(B(2,0)\)，动点\(P\)满足\(k_{PA}\cdot k_{PB}=-1/4\)，求\(P\)的轨迹方程，并验证轨迹的纯粹性与完备性。}
\xiexwei{12mm}

{\ansblockgrayfalse
\begin{ansblock}[2章-练-课时10-中2]
% ans:2章-练-课时10-中2
\ansitem{11}{x²/4+y²=1（x≠±2）}
\ansnote{详解}{设\(P(x,y)\)．斜率存在要求\(x\neq\pm 2\)；\(y^{2}/(x^{2}-4)=-1/4\Rightarrow 4y^{2}=4-x^{2}\)，即\(x^{2}/4+y^{2}=1\)（\(x\neq\pm 2\)）．纯粹性：轨迹上任一点的坐标推导步步可逆，都满足斜率方程（成立）；完备性：满足题设的点都满足\(x^{2}/4+y^{2}=1\)且\(x\neq\pm 2\)，都在轨迹上（成立）（\(y=0\)时斜率积为\(0\neq -1/4\)，端点自然排除）．故轨迹为椭圆\(x^{2}/4+y^{2}=1\)除去长轴两端点\((\pm 2,0)\)．}
\end{ansblock}}

\xiaojie{①识别：给出几何条件（距离、比值、斜率积、垂直平分等）求轨迹方程，判归本题型.}
\bindp {\kaishu ②操作：设动点坐标→把几何条件直译成坐标等式→同解化简→配方（或整理）落标准形；斜率类先写斜率存在条件.}
\bindp {\kaishu ③收束：剔点看三处——分母为零、与定点重合、三点共线退化；求完回代一两个特殊点验算，双向检验纯粹性与完备性.}


% ---------- 探究点三 由方程研究曲线 ----------
\tjdnr{三}{由方程研究曲线}{例\textbf{12}}{中档(知识点三)}{已知曲线\(C\)：\(x^{2}+y^{2}-2|x|=0\)．(1)指出\(C\)的形状；(2)写出\(C\)的对称性、范围、与坐标轴的交点，并求\(C\)围成的区域的面积．}
\xiexwei{12mm}

{\ansblockgrayfalse
\begin{ansblock}[2章-练-课时10-中1]
% ans:2章-练-课时10-中1
\ansitem{12}{(1)两个外切于原点的单位圆（「8」字形）；(2)关于x轴、y轴、原点均对称；x∈[−2,2]，y∈[−1,1]；与y轴仅交于(0,0)，与x轴交于(±2,0)与(0,0)；S=2π}
\ansnote{详解}{(1)按\(|x|\)分段：\(x\geqslant 0\)时\((x-1)^{2}+y^{2}=1\)，\(x<0\)时\((x+1)^{2}+y^{2}=1\)，两圆圆心\((\pm 1,0)\)、半径1，外切于原点．(2)\(|x|\)使\(C\)关于\(y\)轴对称；两圆连心线在\(x\)轴上，故关于\(x\)轴对称；由两者得关于原点对称．范围：\(x\in[-2,2]\)，\(y\in[-1,1]\)．交点：\(x=0\Rightarrow y^{2}=0\)，仅\((0,0)\)；\(y=0\Rightarrow|x|=2\)或\(0\)，得\((\pm 2,0)\)与\((0,0)\)．两圆仅切于一点无重叠，\(S=2\times\pi\times 1^{2}=2\pi\)．}
\end{ansblock}}

\liB{变式\textbf{13}}{中档(知识点三)}{}{已知圆\(C_1\)：\(x^{2}+(y-m)^{2}=9\)与圆\(C_2\)：\(x^{2}+(y+2)^{2}=1\)有公共点，求实数\(m\)的取值范围．}
\xiexwei{10mm}

\begin{ansblock}[2章-练-课时10-中3]
% ans:2章-练-课时10-中3
\ansitem{13}{m∈[−6,−4]∪[0,2]}
\ansnote{详解}{圆心\(C_1(0,m)\)、\(C_2(0,-2)\)，半径\(r_1=3\)、\(r_2=1\)，圆心距\(|C_1C_2|=|m+2|\)．两圆有公共点当且仅当\(|r_1-r_2|\leqslant|C_1C_2|\leqslant r_1+r_2\)，即\(2\leqslant|m+2|\leqslant 4\)，解得\(2\leqslant m+2\leqslant 4\)或\(-4\leqslant m+2\leqslant -2\)，即\(m\in[-6,-4]\cup[0,2]\)．}
\end{ansblock}

\liB{变式\textbf{14}}{中档(知识点三)}{}{就实数\(k\)的取值，讨论方程\(x^{2}+2y^{2}-2x-8y+k=0\)表示的曲线形状．}
\xiexwei{10mm}

\begin{ansblock}[2章-练-课时10-中4]
% ans:2章-练-课时10-中4
\ansitem{14}{k<9：椭圆；k=9：单点(1,2)；k>9：不表示任何曲线}
\ansnote{详解}{配方：\((x-1)^{2}+2(y-2)^{2}=9-k\)．\(k<9\)时右边为正，方程化为\((x-1)^{2}/(9-k)+(y-2)^{2}/\bigl((9-k)/2\bigr)=1\)，表示椭圆；\(k=9\)时方程为\((x-1)^{2}+2(y-2)^{2}=0\)，表示单点\((1,2)\)（退化椭圆）；\(k>9\)时左边非负而右边为负，无实数解，不表示任何曲线．}
\end{ansblock}

\liB{变式\textbf{15}}{难(知识点三)}{}{双纽线，也称伯努利双纽线，其描述首见于1694年，雅各布·伯努利将其作为椭圆的一种类比来处理．椭圆是由到两个定点距离之和为定值的点的轨迹，而卡西尼卵形线则是由到两定点距离之乘积为定值的点的轨迹，当此定值使得轨迹经过两定点的中点时，轨迹便为伯努利双纽线．伯努利将这种曲线称为lemniscate，为拉丁文中“悬挂的丝带”之意．双纽线在数学曲线领域的地位占有至关重要的地位．双纽线像数字“8”，不仅体现了数学的对称、和谐、简洁、统一的美，同时也具有特殊的有价值的艺术美，是形成其它一些常见的漂亮图案的基石，也是许多设计者设计作品的主要几何元素．曲线\(C\)：\((x^{2}+y^{2})^{2}=4(x^{2}-y^{2})\)是双纽线，则下列结论正确的是\nobreak(\kongwei)〔图注：“8”字形双纽线示意，照录自源〕}

\lxopt{A．曲线\(C\)经过5个整点（横、纵坐标均为整数的点）}

\lxopt{B．曲线\(C\)上任意一点到坐标原点\(O\)的距离都不超过2}

\lxopt{C．曲线\(C\)关于直线\(y=x\)对称的曲线方程为\((x^{2}+y^{2})^{2}=4(y^{2}-x^{2})\)}

\lxopt{D．若直线\(y=kx\)与曲线\(C\)只有一个交点，则实数\(k\)的取值范围为\((-\infty,-1]\cup[1,+\infty)\)}

{\ansblockgrayfalse
\begin{ansblock}[2章-练-课时10-难1]
% ans:2章-练-课时10-难1
\ansitem{15}{BCD}
\ansnote{详解}{A：令\(y=0\)得\(x^{4}=4x^{2}\)，\(x=0\)或\(\pm 2\)，整点\((2,0)\)、\((-2,0)\)、\((0,0)\)三个；若整点满足\(y\neq 0\)，则由B项结论\(x^{2}+y^{2}\leqslant 4\)只能\(y=\pm 1\)、\(x\in\{0,\pm 1,\pm 2\}\)：代入\((x^{2}+1)^{2}=4(x^{2}-1)\)，记\(X=x^{2}\)得\(X^{2}-2X+5=0\)，判别式\(4-20<0\)无解——\(y\neq 0\)无整点．故曲线仅过3个整点，“5个整点”错．B：由\((x^{2}+y^{2})^{2}=4(x^{2}-y^{2})\leqslant 4(x^{2}+y^{2})\)得\(x^{2}+y^{2}\leqslant 4\)，即\(|OP|\leqslant 2\)（成立）．C：方程中\(x\)、\(y\)互换得\((x^{2}+y^{2})^{2}=4(y^{2}-x^{2})\)，即关于\(y=x\)对称的曲线（成立）．D：\(y=kx\)代入：\(x^{4}(1+k^{2})^{2}=4x^{2}(1-k^{2})\)．\(x=0\)恒给交点\((0,0)\)；\(|x|>0\)的解须\(x^{2}=4(1-k^{2})/(1+k^{2})^{2}>0\)，当且仅当\(|k|<1\)．故\(|k|\geqslant 1\)时只有原点一个交点；\(|k|<1\)时另有\(\pm\)两交点（\((0,0)\)之外），不止一个（成立）（\(k=\pm 1\)时\(x^{4}\cdot 2=0\)仅原点，含入界点正确）．故选：BCD．}
\end{ansblock}}

\liB{变式\textbf{16}}{难(知识点三)}{}{关于曲线\(C\)：\((x-m)^{2}+(y-2m)^{2}=n^{2}/2\)有以下五个结论：}
\bindp ①当\(m=1\)时，曲线\(C\)表示圆心为\((1,2)\)，半径为\((\sqrt{2}/2)|n|\)的圆；

\bindp ②当\(m=0\)，\(n=2\)时，过点\((3,3)\)向曲线\(C\)作切线，切点为\(A\)，\(B\)，则直线\(AB\)的方程为\(3x+3y-2=0\)；

\bindp ③当\(m=1\)，\(n=\sqrt{2}\)时，过点\((2,0)\)向曲线\(C\)作切线，则切线方程为\(y=-(3/4)(x-2)\)；

\bindp ④当\(n=m\neq 0\)时，曲线\(C\)表示圆心在直线\(y=2x\)上的圆系，且这些圆的公切线方程为\(y=x\)或\(y=7x\)；

\bindp ⑤当\(n=4\)，\(m=0\)时，直线\(kx-y+1-2k=0\)（\(k\in R\)）与曲线\(C\)表示的圆相离．

\bindp 以上正确结论的序号为\kongbai{}．

{\ansblockgrayfalse
\begin{ansblock}[2章-练-课时10-难2]
% ans:2章-练-课时10-难2
\ansitem{16}{②④}
\ansnote{详解}{①\(m=1\)时\((x-1)^{2}+(y-2)^{2}=n^{2}/2\)，但\(n=0\)时退化为点\((1,2)\)，不（恒）表示圆——以偏概全，错误．②\(m=0\)、\(n=2\)：\(C\)：\(x^{2}+y^{2}=2\)．记\(D(3,3)\)，\(|DO|=3\sqrt{2}\)，切线长\(|DA|=|DB|=\sqrt{18-2}=4\)，切点\(A\)、\(B\)即\(C\)与圆\((x-3)^{2}+(y-3)^{2}=16\)（圆心\(D\)半径4）的公共点，公共弦方程由两圆方程相减：\(x^{2}+y^{2}-2-(x^{2}+y^{2}-6x-6y+2)=0\)，即\(3x+3y-2=0\)，正确．③\(m=1\)、\(n=\sqrt{2}\)：\(C\)：\((x-1)^{2}+(y-2)^{2}=1\)．过\((2,0)\)的切线：竖直直线\(x=2\)到圆心\((1,2)\)距离为\(1=\)半径，已是切线；另设\(y=k(x-2)\)：\(|k-2-2k|/\sqrt{k^{2}+1}=1\Rightarrow(k+2)^{2}=k^{2}+1\Rightarrow k=-3/4\)，即\(3x+4y-6=0\)也是切线．结论只给后者、漏\(x=2\)，“切线方程为”不完备，错误．④\(n=m\neq 0\)：\((x-m)^{2}+(y-2m)^{2}=m^{2}/2\)，圆心\((m,2m)\)恒在\(y=2x\)上，半径\((\sqrt{2}/2)|m|\)．圆心到\(y=x\)距离\(=|m-2m|/\sqrt{2}=(\sqrt{2}/2)|m|=\)半径；到\(y=7x\)距离\(=|7m-2m|/\sqrt{50}=(\sqrt{2}/2)|m|=\)半径——两直线都是圆系公切线．穷尽性补核（源未给）：设公切线\(y=kx+b\)对一切\(m\)成立，则\([m(k-2)+b]^{2}=\frac{1}{2}m^{2}(k^{2}+1)\)恒成立，比较系数：\(b=0\)且\((k-2)^{2}=(k^{2}+1)/2\Rightarrow k^{2}-8k+7=0\Rightarrow k=1\)或\(7\)；竖直线\(x=c\)代入\(|m-c|=(\sqrt{2}/2)|m|\)不能对一切\(m\)成立，无．故公切线恰为\(y=x\)、\(y=7x\)，正确．⑤\(n=4\)、\(m=0\)：\(C\)：\(x^{2}+y^{2}=8\)，直线\(kx-y+1-2k=0\)即\(k(x-2)-(y-1)=0\)恒过定点\((2,1)\)，而\(4+1=5<8\)，定点在圆内，直线恒与圆相交而非相离，错误．综上，正确的是②④．}
\end{ansblock}}

\xiaojie{①识别：给出方程研究形状、对称、范围、交点，或含参方程、曲线系方程定形分类，判归本题型.}
\bindp {\kaishu ②操作：先定定义域再化简；含绝对值分段去绝对值；含参方程配方读系数定形；两圆公共点比圆心距与半径和、差；切点弦用过定点的圆系公共弦（两方程相减）.}
\bindp {\kaishu ③收束：含参定形按右边正负分档不漏退化档（单点）；多结论逐项独立验证，“以偏概全”与“漏解”（如漏竖直切线）都判错.}

% ============================================================
% 课堂评价 G5（恒5＝3单选+2填空＝导槽 G1~G5；ketang 新制顺排可跨栏，禁 \vbox 装栏）
% 印面号 17—21；多选门：本片正文多选 0（10-难1 槽型「多结论选择」不计多选门，答案 BCD）
% ============================================================
\begin{ketang}{本组共5题（第17—21题），17—19为单选，20—21为填空．}

\jiancestem{{\fontsize{11.4pt}{14pt}\selectfont\heihao {\numboldjian 17}．}\hspace{4.1mm}\tieside{中档(知识点二)}\hspace{2.2mm}已知点\(P\)是圆\(O\)：\(x^{2}+y^{2}=4\)上的动点，作\(PH\perp y\)轴于点\(H\)，则线段\(PH\)的中点\(M\)的轨迹方程为\nobreak(\kongwei)}

\bindopt {\fontsize{10.5pt}{\qpTJgLeadPingjia}\selectfont \makebox[\dimexpr0.5\linewidth-1em\relax][l]{A．\(x^{2}/4+y^{2}=1\)}\makebox[\dimexpr0.5\linewidth-1em\relax][l]{B．\(x^{2}/16+y^{2}=1\)}\par}

\bindopt {\fontsize{10.5pt}{\qpTJgLeadPingjia}\selectfont \makebox[\dimexpr0.5\linewidth-1em\relax][l]{C．\(x^{2}+y^{2}/16=1\)}\makebox[\dimexpr0.5\linewidth-1em\relax][l]{D．\(x^{2}+y^{2}/4=1\)}\par}

\begin{ansblock}[2章-导-课时10-G1]
% ans:2章-导-课时10-G1
\ansitem{17}{D}
\ansnote{详解}{设\(M(x,y)\)，\(P(x_0,y_0)\)，则\(x_0^{2}+y_0^{2}=4\)．由\(PH\perp y\)轴知\(H(0,y_0)\)；\(M\)为\(PH\)中点，得\(x=x_0/2\)、\(y=y_0\)，即\(x_0=2x\)、\(y_0=y\)．代入圆方程：\(4x^{2}+y^{2}=4\)，整理得\(x^{2}+y^{2}/4=1\)．故选：D．}
\end{ansblock}

\jiancestem{{\fontsize{11.4pt}{14pt}\selectfont\heihao {\numboldjian 18}．}\hspace{4.1mm}\tieside{中档(知识点二)}\hspace{2.2mm}已知圆\(C\)：\((x-2)^{2}+y^{2}=64\)，\(F(-2,0)\)为圆内一点，将圆折起使得圆周过点\(F\)（如图），然后将纸片展开，得到一条折痕\(l\)，这样继续下去将会得到若干折痕，观察这些折痕围成的轮廓是一条圆锥曲线，则该圆锥曲线的方程为\nobreak(\kongwei)〔图注：折叠示意图，照录自源〕}

\zhuzhu{定义提示卡：椭圆的定义——平面内与两个定点\(F_1\)、\(F_2\)的距离之和等于常数（大于\(|F_1F_2|\)）的点的轨迹；折痕恰是折叠对应点连线的垂直平分线．本题主解将用到（2.5.1 系统学习）．}

\bindopt {\fontsize{10.5pt}{\qpTJgLeadPingjia}\selectfont \makebox[\dimexpr0.5\linewidth-1em\relax][l]{A．\(x^{2}/16+y^{2}/12=1\)}\makebox[\dimexpr0.5\linewidth-1em\relax][l]{B．\(x^{2}/4+y^{2}=1\)}\par}

\bindopt {\fontsize{10.5pt}{\qpTJgLeadPingjia}\selectfont \makebox[\dimexpr0.5\linewidth-1em\relax][l]{C．\(x^{2}/4+y^{2}/3=1\)}\makebox[\dimexpr0.5\linewidth-1em\relax][l]{D．\(x^{2}/16+y^{2}/4=1\)}\par}

\begin{ansblock}[2章-导-课时10-G2]
% ans:2章-导-课时10-G2
\ansitem{18}{A}
\ansnote{详解}{\(F(-2,0)\)、\(C(2,0)\)，记\(F\)关于折痕\(l\)的对称点为\(A\)，\(l\)与\(AC\)交于点\(P\)，则\(A\)在圆周上，\(l\)为\(AF\)的垂直平分线，故\(|PA|=|PF|\)．于是\(|PF|+|PC|=|PA|+|PC|=|AC|=8>|FC|=4\)，点\(P\)的轨迹是以\(F\)、\(C\)为左、右焦点的椭圆，\(2a=8\)、\(2c=4\)，\(a=4\)、\(c=2\)、\(b=2\sqrt{3}\)，方程为\(x^{2}/16+y^{2}/12=1\)．故选：A．}
\end{ansblock}

\jiancestem{{\fontsize{11.4pt}{14pt}\selectfont\heihao {\numboldjian 19}．}\hspace{4.1mm}\tieside{中档(知识点三)}\hspace{2.2mm}已知曲线\(C\)：\(\sqrt{x-1}\cdot\ln(x^{2}+y^{2}-1)=0\)，直线\(y=kx-1\)与曲线\(C\)恰有两个交点，则\(k\)的取值集合为\nobreak(\kongwei)}

\bindopt {\fontsize{10.5pt}{\qpTJgLeadPingjia}\selectfont \makebox[\dimexpr0.5\linewidth-1em\relax][l]{A．\(\{0,1\}\)}\makebox[\dimexpr0.5\linewidth-1em\relax][l]{B．\(\{0,2\}\)}\par}

\bindopt {\fontsize{10.5pt}{\qpTJgLeadPingjia}\selectfont \makebox[\dimexpr0.5\linewidth-1em\relax][l]{C．\((0,2)\)}\makebox[\dimexpr0.5\linewidth-1em\relax][l]{D．\((0,1)\cup(1,2)\)}\par}

{\ansblockgrayfalse
\begin{ansblock}[2章-导-课时10-G3]
% ans:2章-导-课时10-G3
\ansitem{19}{D}
\ansnote{详解}{曲线定义域：\(x-1\geqslant 0\)且\(x^{2}+y^{2}-1>0\)．乘积为零当且仅当\(\sqrt{x-1}=0\)或\(\ln(x^{2}+y^{2}-1)=0\)．前者给\(x=1\)（此时\(x^{2}+y^{2}-1=y^{2}>0\)须\(y\neq 0\)），即直线\(x=1\)除去\((1,0)\)；后者给\(x^{2}+y^{2}=2\)且\(x\geqslant 1\)，即圆\(x^{2}+y^{2}=2\)的右半弧（含端点\((1,\pm 1)\)，该两点同时落在\(x=1\)上）．直线\(y=kx-1\)恒过定点\((0,-1)\)：与\(x=1\)（\(y\neq 0\)）交于\((1,k-1)\)，\(k\neq 1\)时是曲线上一交点；又\((0,-1)\)到原点距离\(1<\sqrt{2}\)，直线与整圆恒交于两点，其中落在右半弧上的：临界为直线过弧端点\((1,-1)\)（\(k=0\)）与\((1,1)\)（\(k=2\)），端点处弧上点与\(x=1\)上点重合、两交并一，故须\(0<k<2\)（严格）；再排除\(k=1\)（此时\((1,k-1)=(1,0)\)不在曲线上，曲线只剩弧上一个交点，\(k=1\)时直线与圆\(x\geqslant 1\)半弧恰一交点）．综上\(0<k<2\)且\(k\neq 1\)，故选：D．}
\end{ansblock}}

\jiancestem{{\fontsize{11.4pt}{14pt}\selectfont\heihao {\numboldjian 20}．}\hspace{4.1mm}\tieside{中档(知识点二)}\hspace{2.2mm}已知等腰三角形\(ABC\)的顶点为\(A(4,2)\)，底边的一个端点为\(B(5,3)\)，则底边的另一个端点\(C\)的轨迹方程为\kongbai{}．}
\xiexwei{10mm}

{\ansblockgrayfalse
\begin{ansblock}[2章-导-课时10-G4]
% ans:2章-导-课时10-G4
\ansitem{20}{x²+y²−8x−4y+18=0（除去点(3,1)与(5,3)）}
\ansnote{详解}{设\(C(x,y)\)．由\(|AC|=|AB|\)得\(\sqrt{(x-4)^{2}+(y-2)^{2}}=\sqrt{(4-5)^{2}+(2-3)^{2}}=\sqrt{2}\)，平方化简得\(x^{2}+y^{2}-8x-4y+18=0\)（圆心\((4,2)\)、半径\(\sqrt{2}\)的圆）．剔除：\(C\)与\(B\)重合，去\((5,3)\)；\(A\)、\(B\)、\(C\)三点共线时不构成三角形——\(k_{AB}=(3-2)/(5-4)=1\)，直线\(AB\)：\(y-2=1\cdot(x-4)\)即\(x-y-2=0\)，与圆联立解得交点\((3,1)\)与\((5,3)\)，故须除去\((3,1)\)、\((5,3)\)．轨迹方程为\(x^{2}+y^{2}-8x-4y+18=0\)（\(x-y-2\neq 0\)，即除去\((3,1)\)、\((5,3)\)）．}
\end{ansblock}}

\jiancestem{{\fontsize{11.4pt}{14pt}\selectfont\heihao {\numboldjian 21}．}\hspace{4.1mm}\tieside{简单(知识点二)}\hspace{2.2mm}求到\(y\)轴的距离等于4的点的轨迹的方程．}
\xiexwei{10mm}

\begin{ansblock}[2章-导-课时10-G5]
% ans:2章-导-课时10-G5
\ansitem{21}{x=4 或 x=−4（即 x=±4，两条平行于y轴的直线）}
\ansnote{详解}{设动点\(M(x,y)\)．\(M\)到\(y\)轴距离为\(|x|\)，条件\(|x|=4\)即\(x=4\)或\(x=-4\)．检验两性：直线\(x=\pm 4\)上任意一点到\(y\)轴距离均为4（纯粹性成立）；到\(y\)轴距离为4的点横坐标为\(\pm 4\)，都在两直线上（完备性成立）．故轨迹方程为\(x=\pm 4\)．}
\end{ansblock}

\end{ketang}

% ---- 尾块（\tailfill 置 \end{multicols} 前；无参＝内容尾随框，件侧不擅自定高） ----
\tailfill

\end{multicols}
\end{document}
